\documentstyle[12pt,newsymb]{article}
\newcommand\eref[1]{(\ref{#1})}
\def\a{\alpha}
\def\d{\delta}
\def\i{\infty}
\def\t{\theta}

\begin{document}

\title{On a Coefficient Conjecture of Brannan}
\author{Roger W.\ Barnard, Kent Pearce, and William Wheeler}
\date{}
\maketitle

\begin{abstract}
In 1972, D.A.\ Brannan conjectured that all of the odd coefficients,
$a_{2n+1}$, of the power series $(1+xz)^\a/(1-z)$ were dominated by
those of the series $(1+z)^\a/(1-z)$ for the parameter range
$0<\a<1$, after having shown that this was not true for the even
coefficients.  He verified the case when $2n+1=3$.  The case when
$2n+1=5$ was verified in the mid-eighties by J.G.\ Milcetich.  In
this paper, we verify the case when $2n+1=7$ using classical Sturm
sequence arguments and some computer algebra.
\end{abstract}

1980 AMS Subject Classification: 30C50\\

Keywords: coefficient conjecture, computer algebra

\newpage

\baselineskip=18pt

\noindent{\bf Introduction.}

For $k\geq 2$ let $V_k$ denote the class of locally univalent
analytic functions
\begin{equation}
\label{1}
f(z)=z+a_2z^2+a_3z^3+\cdots
\end{equation}
which map $|z|<1$ conformally onto a domain whose boundary rotation is
at most $k\pi$.  (See [Pa] for the definition and basic properties of
the class $V_k$.)

The function
$$
f_k(z)=\frac{1}{k}\left[\left(\frac{1+z}{1-z}\right)^{\frac{k}{2}}
-1\right]=\sum^\i_{n=1}A_nz^n
$$
belongs to $V_k$.  The coefficient conjecture for the class $V_k$ was
that for a function (1) in $V_k$ that
\begin{equation}
\label{2}
|a_n|\leq A_n,\;\; (n\geq 1).
\end{equation}
This conjecture was verified for $n=2$ by Pick (see [Le]), for $n=3$
by Lehto [Le] in 1952 and for $n=4$ by Schiffer and Tammi [ScTa] in
1967, Lonka and Tammi [LoTa] in 1968 and Brannan [Br1] in 1969.

Using extreme point theory arguments, Brannan, Clunie and Kirwan
[BrClKi] showed in 1973 that (2) can be reduced to showing that for
$$
\Phi(\a,x;z)=\left(\frac{1+xz}{1-z}\right)^\a=
\sum^\i_{n=1}B_n(\a,x)z^n
$$
that
\begin{equation}
\label{3}
|B_n(\a,x)|\leq B_n(\a,1),\;\; (n\geq 1)
\end{equation}
for $\a\geq 1$, $|x|=1$.  Brannan, Clunie and Kirwan showed that (3)
holds for $1\leq n\leq 13$, which implies (2) for $2\leq n\leq 14$.

In 1972 Aharonov and Friedland [AhFr] considered a related
coefficient inequality.  Let
$$
\Psi(\a,x;z)=\frac{(1+xz)^\a}{1-z}=\sum^\i_{n=1}A_n(\a,x)z^n.
$$
In [AhFr] it was shown, by a long technical argument, that
\begin{equation}
\label{4}
|A_n(\a,x)|\leq A_n(\a,1),\;\; (n\geq 1)
\end{equation}
for $\a\geq 1,|x|=1$, which implies (3) and, hence, by the work in
[BrClKi], also implies (2).  Later, in 1973 Brannan [Br2] gave a
short, elegant proof that (4) holds for $\a\geq 1$, $|x|=1$.

In [Br2] Brannan also considered the question about whether (4) holds
for $0<\a<1$, $|x|=1$.  He showed there the unexpected result that
for each $\a$, $0<\a<1$, there exists an $n_\a$ such that
\begin{equation}
\label{5}
\max_{|x|=1}{Re}\; A_{2n}(\a,x)>A_{2n}(\a,1)
\end{equation}
for $n>n_\a$, that is, that (4) fails for even coefficients when
$0<\a<1$.

Brannan showed, using an inequality for quadratic trigonometric
polynomials, that
$$
|A_3(\a,x)|\leq A_3(\a,1)
$$
for $0<\a<1$ and he conjectured, based on numerical data, that\\

\noindent{\bf Brannan's Conjecture}
\begin{equation}
\label{6}
|A_{2n+1}(\a,x)|\leq A_{2n+1}(\a,1),\;\; (n\geq 1)
\end{equation}
for $0<\a<1$, $|x|=1$.

Brannan's conjecture has been verified for $n=2$, that is, for
$2n+1=5$, by Milcetich [Mi], who employed a lengthy argument based
on a result of Brown and Hewitt [BrHe] for positive trigonometric
sums.

In this paper, we will establish Brannan's conjecture for $n=3$, that
is, for $2n+1=7$.  The method we will employ is based largely on (i)
a judicious rearrangement of the coefficients $A_n(\a,x)$ over
carefully chosen subintervals of (0,1), the domain of $\a$, (ii) an
application of Sturm sequences to verify the nonnegativity of those
rearrangements and (iii) using a computer algebra program (in this
case Maple) to generate the coefficients $A_n(\a,x)$ and the Sturm
sequences.\\

\noindent{\bf Section 1.}\\

Brannan's coefficient inequality (6) is equivalent to
\begin{equation}
\label{7}
A^2_{2n+1}(\a,1)-|A_{2n+1}(\a,x)|^2\geq 0
\end{equation}
for $0<\a<1$, $|x|=1$.  We will let $F_{2n+1}(\a,x)$ denote the
left-hand side of (7) and we will show for $2{n+1}=7$ that
$F_{2n+1}(\a,x)\geq 0$.

We note that
\begin{eqnarray*}
& \displaystyle{\frac{(1+xz)^\a}{(1-z)}} & =\sum^\i_{n=0} 
\frac{(-\a)_n(-1)^nx^n}{n!}z^n\sum^\i_{n=0}z^n\\[1ex]
&& =\sum^\i_{n=0}\sum^n_{k=0}\frac{(-\a)_k(-1)^kx^k}{k!}z^n\\[1ex]
&& =\sum^\i_{n=0}A_n(\a,x)z^n,
\end{eqnarray*}
where $(a)_k$ denotes the Pockhammer symbol, which is defined as
$$
(a)_k=\left\{ \begin{array}{cl} 1 & k=0\\[1ex] a(a+1)\cdots(a+k-1) &
k>0\end{array}\right. .
$$
Hence, we can write $F_N(\a,x)$ as
\begin{eqnarray*}
F_N(\a,x)\indent & = & \sum^N_{k=0}\; \frac{(-\a)_k(-1)^k}{k!}\;
\sum^N_{k=0}\; \frac{(-\a)_k(-1)^k}{k!}\; -\\[1ex]
&& \sum^N_{k=0}\; \frac{(-\a)_k(-1)^kx^k}{k!}\; \sum^N_{k=0}\; 
\frac{(-\a)_k(-1)^k\bar{x}^k}{k!}
\end{eqnarray*}

\newpage

\begin{eqnarray*}
& = & \sum^{2*N}_{k=0}\; \sum^k_{j=0}\; 
\frac{(-\a)_j(-\a)_{k-j}(-1)^k\d_j\d_{k-j}}{k!(k-j)!}\; -\\[1ex]
&& \sum^{2*N}_{k=0}\; \sum^k_{j=0}\;
\frac{(-\a)_j(-\a)_{k-j}(-1)^k\d_j\d_{k-j}x^{2*j-k}}{k!(k-j)!}
\end{eqnarray*}
where
$$
\d_j=\left\{ \begin{array}{lc} 1 & 0\leq j\leq N\\[1ex] 
0 & N+1\leq j\leq 2*N\end{array}\right.
$$
Since $F_N(\a,x)$ is real, we can write, setting $x=e^{i\t}$, 
\begin{eqnarray}
\nonumber
F_N(\a,x) & = & \sum^{2*N}_{k=0}\; \sum^k_{j=0}\; 
\frac{(-\a)_j(-\a)_{k-j}(-1)^k\d_j\d_{k-j}}{k!(k-j)!}\; -\\[1ex]
\label{8}
&& \sum^{2*N}_{k=0}\; \sum^k_{j=0}\;
\frac{(-\a)_j(-\a)_{k-j}(-1)^k\d_j\d_{k-j}\cos((2*j-k)\t)}{k!(k-j)!}
\end{eqnarray}

The following Maple Procedure can be used to generate the
coefficients of $F_N(\a,x)$, where $x=e^{i\t}$,
\vspace{8pt}

\begin{center} Procedure 1 \end{center}
\vspace{8pt}

{\sf \begin{tabular}{l}
F:=proc(N)\\
local i, j, a, csum, dsum, temp;\\
global c;\\
\indent a[0]:=1;\\
\indent for i from 1 to N do a[i]:=a[i-1]*(-alpha+i-1)*(-1)/i od;\\
\indent for i from N+1 to 2*N do a[i]:=0 od;\\
\indent csum:=0; dsum:=0;\\
\indent for i from 0 to N do csum:= csum+a[i] od;\\
\end{tabular}}

{\sf \begin{tabular}{l}
\indent for i from 0 to 2*N do\\
\indent for j from 0 to i do\\
\indent \indent dsum:= dsum + a[j]*a[i-j]* cos ((2*j-i)*theta);\\
\indent od;\\
\indent od;\\
\indent temp:= collect(csum*csum-dsum,alpha);\\
\indent for i from 0 to (2*N-1) do c[i]:= coeff(temp,alpha,i) od;\\
\indent temp;\\
end;\\
\end{tabular}}\\

Using Procedure 1, we obtain for $N=7$ that $F_7(\a,x)=
\displaystyle{\sum^{13}_{k=1}}c_k(\t)\a^k$ where each $c_k(\t)$ is a
trigonometric polynomial of the form $c_k(\t)=\displaystyle{
\sum^7_{j=0}}a_{kj}\cos(j\t)$ with rational coefficients $a_{kj}$.
We will show that $F_7(\a,x)\geq 0$ for $0<\a<1$ by subdividing the
domain of $\a$ into subintervals $0<\a\leq t_0$ and $t_0<\a<1$, where
$t_0=2/5$.  We will show that $F_7(\a,x)\geq 0$ on each subinterval.

First for the case $0<\a\leq t_0$ we will show the following:
$$ 
c_1(\t)\geq 0,\;\; \frac{7}{10} c_1(\t)+c_2(\t) t_0\geq 0,
\eqno{(9.1)}
$$
$$
\frac{1}{10} c_1(\t)+c_3(\t)t^2_0\geq 0,\;\; \frac{1}{10}
c_1(\t)+c_3(\t)t^2_0+c_4(\t)t^3_0\geq 0, \eqno{(9.2)}
$$
$$
\frac{1}{5} c_1(\t)+c_5(\t)t^4_0\geq 0,\;\; \frac{1}{5}
c_1(\t)+c_5(\t)t^4_0+c_6(\t)t^5_0\geq 0, \eqno{(9.3)}
$$
$$
c_7(\t)\geq 0,\;\; c_7(\t)+c_8(\t)t_0\geq 0, \eqno{(9.4)}
$$
$$
c_9(\t)\geq 0,\;\; c_9(\t)+c_{10}(\t)t_0\geq 0, \eqno{(9.5)}
$$
$$
c_{11}(\t)\geq 0,\;\; c_{11}(\t)+c_{12}(\t)t_0\geq 0, \eqno{(9.6)}
$$
$$
c_{13}(\t)\geq 0.  \eqno{(9.7)}
$$
It will follow then that for $0<\a\leq t_0$ we have

\setcounter{equation}{9}

\begin{equation}
\label{10}
\begin{array}{ll} F_7(\a,x) 
& =\left[\frac{7}{10} c_1(\t)+c_2(\t)\a\right]
\a+\left[\frac{1}{10}c_1(\t)+c_3(\t)\a^2+c_4(\t)\a^3 \right]\a\\[1ex]
& \mbox{} + \left[\frac{1}{5}c_1(\t)+c_5(\t)\a^4+c_6(\t)\a^5\right]
\a+\left[c_7(\t)+c_8(\t)\a\right]\a^7\\[1ex] 
& \mbox{} + \left[c_9(\t)+c_{10}(\t)\a\right]\a^9+
\left[c_{11}(\t)+c_{12}(\t)\a\right]\a^{11}\\[1ex] 
& \mbox{} + c_{13}(\t)\a^{13}\geq 0\end{array}
\end{equation}

The inequalities (9) imply (10) because they imply that each of the
terms in brackets in (10) are non-negative.  The non-negativity of
the bracketed terms of the form $[c_i(\t)+c_{i+1}(\t)\a]$ follows
from (9.1), (9.4), (9.5) and (9.6) because the terms are linear in
$\a$ and, hence they take their minimum at either $\a=0$ or else at
$\a=t_0$.

Since $[c_3(\t)+c_4(\t)\a]$ is linear in $\a$, it takes its minimum
at either $\a=0$ or else at $\a=t_0$.  Thus, we have
\begin{equation}
\label{11}
\frac{1}{10} c_1(\t) + (c_3(\t)+c_4(\t)\a) \a^2 \geq \frac{1}{10}
c_1(\t) + \min_{0<s\leq t_0} \{c_3(\t)+c_4(\t)s\} \a^2.
\end{equation}
The right-hand side of (11) is linear in $\a^2$, and hence takes its
minimum at either $\a=0$ or else at $\a=t_0$.  Therefore, the
right-hand side of (11) is non-negative by (9.2).

Similarly, since $[c_5(\t)+c_6(\t)\a]$ is linear in $\a$, it takes
its minimum at either $\a=0$ or else at $\a=t_0$.  Thus, we have 
\begin{equation}
\label{12}
\frac{1}{5} c_1(\t) + (c_5(\t)+c_6(\t)\a) \a^4 \geq \frac{1}{5}
c_1(\t) + \min_{0<s\leq t_0} \{c_5(\t)+c_6(\t)s\} \a^4.
\end{equation}
The right-hand side of (12) is linear in $\a^4$, and hence takes its
minimum at either $\a=0$ or else at $\a=t_0$.  Therefore, the
right-hand side of (12) is non-negative by (9.3).

Thus, to complete the case $0<\a\leq t_0$ we will need to establish
(9).  We will transform each of the trigonometric coefficients
$c_i(\t)$, which are polynomials in $\cos(n\t)$, to polynomials in
$\cos\t$ and then by a change of variable to polynomials $e_i(x)$,
$-1\leq x\leq 1$.  To verify the non-negativity of the linear
combinations of trigonometric coefficients $c_i(\t)$ specified in
(9), we will establish the non-negativity of the same linear
combinations of polynomials $e_i(x)$.

The following two Maple procedures can be used to: (i) transform the
trigonometric coefficients $c_i(\t)$ to the polynomials $e_i(x)$; and
(ii) compute the number of roots of a polynomial {\sf p} on the interval
$(-1,1]$ via a Sturm sequence argument.
\vspace{8pt}

\begin{center} Procedure 2 \end{center}
\vspace{8pt}

{\sf \begin{tabular}{l}
G:=proc(N)\\
local i, t, temp;\\
global c, e;\\
\indent for i from 0 to 2*N - 1 do\\
\indent\indent t[i] := expand(c[i]);\\
\indent\indent e[i] := subs(cos(theta) = x, t[i]);\\
\indent od;\\
\indent temp;\\
end;\\
\end{tabular}}\\

\noindent Here {\sf N} is chosen the same as in Procedure 1.\\

\begin{center} Procedure 3 \end{center}
\vspace{8pt}

{\sf \begin{tabular}{l}
H := proc(p)\\
local s;\\
global lc, nr;\\
\indent lc := roots(p, x);\\
\indent s := sturmseq(p, x);\\
\indent nr := sturm(s, x, -1, 1)\\
end;\\
\end{tabular}}\\

\noindent The library call {\em readlib(sturm)} must be loaded prior
to applying the procedure.\\

If the polynomial $e(x)$, created from linear combinations of the
$e_i(x)$ after applying Procedure 2, is assigned to the variable {\sf
p}, then Procedure 3 will compute both the number of roots of $e(x)$
on the interval $(-1,1]$ and the location of the rational roots of
$e(x)$.  We will see that the conclusion of this application of
Procedure 3 is that the polynomial $e(x)$ is non-negative on $[-1,1]$
with $e(x)=0$ only for $x=1$.  This check can be confirmed for each
of the polynomials $e(x)$ which arise as linear combinations of the
polynomials $e_i(x)$, where the linear combinations are specified as
in (9), and, thus, complete the case $0<\a\leq t_0$.

To illustrate the utility of using computer algebra software to
establish the inequalities in (9) we will explicitly demonstrate the
process for the inequality (9.1).  Applying Procedure 1 to compute
the trigonometric coefficients $c_i(\t)$ of $F_7(\a,x)$, we obtain
\begin{eqnarray*}
& c_1(\t) & = - \frac{2}{5} \cos(5\t) + \cos(2 \t) - \frac{2}{3}
\cos(3\t) + \frac{1}{3} \cos(6\t)\\[1ex]
&& \mbox{} + \frac{1}{2} \cos(4\t) - 2\cos(\t) - \frac{2}{7} 
\cos(7\t) + \frac{319}{210}
\end{eqnarray*}
From Procedure 2 we obtain
\begin{eqnarray*}
& e(x) & = e_1(x) = \frac{128}{5}x^5-\frac{32}{3}x^3+4\;x^2+ 
\frac{24}{35}\\[1ex]
&& \mbox{} + \frac{32}{3}x^6-12\;x^4-\frac{128}{7}x^7.
\end{eqnarray*}
Then, Procedure 3 yields
$$
\begin{array}{lccc} >{\sf H(p)};\\[1ex] &&& 1\\[1ex] >{\sf lc}\\[1ex]
&&& [[1,1]]\end{array}
$$
The value returned by the procedure call {\sf H(p)} is the number of
roots of {\sf p} on $(-1,1]$ and the value returned by {\sf lc} is
the interval location of the rational roots of {\sf p}.

For the second half of (9.1), we have
\begin{eqnarray*}
& \frac{7}{10}c_1(\t)+c_2(\t)t_0 & = \frac{128}{525} \cos(5\t) - 
\frac{37}{210} \cos(2\t) + \frac{451}{1575} \cos(3\t)\\[1ex]
&& \mbox{}-\frac{292}{1575}\cos(60)-\frac{197}{700}\cos(4\t)- 
\frac{5}{7}\cos(\t)+\frac{2}{25}\cos(7\t)+\frac{523}{700}
\end{eqnarray*}
From Procedure 2 we obtain
\begin{eqnarray*}
& e(x)=\frac{7}{10}e_1(x)+e_2(x)t_0 & =-\frac{2656}{525}x^5+ 
\frac{236}{315}x^3-\frac{151}{105}x^2\\[1ex]
&& \mbox{} +\frac{1303}{1575}-\frac{9344}{1575}x^6+ 
\frac{698}{105}x^4+\frac{128}{25}x^7-\frac{32}{35}x
\end{eqnarray*}
Then, Procedure 3 yields
$$
\begin{array}{lccc} >{\sf H(p)};\\[1ex] &&& 1;\\[1ex] 
>{\sf lc}\\[1ex] &&& [[1,1]]\end{array}
$$

We have that each linear combination $e(x)$ has only one root on
$(-1,1]$ and that root is at $x=1$.  Since we can explicitly observe
that each $e(0)>0$, we can conclude that each $e(x)$ is non-negative
on $[-1,1]$.  Therefore, we have that (9.1) holds.

For the case $t_0<\a<1$ we make the substitution $\a=\beta+t_0$.
Then, we have $F_7(\a,x)=G_7(\beta,x)=\sum^{13}_{k=0}d_k(\t)\beta^k$
where each $d_k(\t)$ is a trigonometric polynomial of the form
$d_k(\t)=\sum^7_{j=0}b_{kj}\cos(j\t)$ and $0<\beta<t_1=3/5$.  It
will suffice to show for this case that
$$ 
d_0(\t)\geq 0,\;\;d_0(\t)+d_2(\t)t^2_1\geq 0 \eqno{(13.1)}
$$
$$
d_1(\t)\geq 0,\;\;\frac{2}{3}d_1(\t)+d_5(\t)t^4_1\geq 0\eqno{(13.2)}
$$
$$
\frac{1}{3}d_1(\t)+d_6(\t)t^5_1\geq 0,\;\;\frac{1}{3}d_1(\t)+ 
d_6(\t)t^5_1+d_8(\t)t^7_1\geq 0 \eqno{(13.3)}
$$
$$
d_3(\t)\geq 0,d_3(\t)+d_4(\t)t_1\geq 0 \eqno{(13.4)}
$$
$$
d_7(\t)\geq 0 \eqno{(13.5)}
$$
$$
d_9(\t)\geq 0,\;\; d_9(\t)+d_{10}(\t)t_1\geq 0, \eqno{(13.6)}
$$
$$
d_{11}(\t)\geq 0,\;\; d_{11}(\t)+d_{12}(\t)t_1\geq 0, \eqno{(13.7)}
$$
$$
d_{13}(\t)\geq 0.  \eqno{(13.8)}
$$
For then, it will follow that for $0<\beta<t_1$ we have

\setcounter{equation}{13}

\begin{equation}
\label{14}
\begin{array}{ll} G_7(\beta,x) 
& =\left[d_0(\t)+d_2(\t)\beta^2\right] + 
\left[\frac{2}{3}d_1(\t)+d_5(\t)\beta^4\right]\beta\\[1ex]
& \mbox{} + \left[\frac{1}{3}d_1(\t)+d_6(\t)\beta^5+d_8(\t)
\beta^7\right]\beta+\left[d_3(t)+d_4(\t)\beta\right]\beta^3\\[1ex]
& \mbox{} + d_7(\t)\beta^7+\left[d_9(\t)+d_{10}(\t)\beta\right]
\beta^9+\left[d_{11}(\t)+d_{12}(\t)\beta\right]\beta^{11}\\[1ex]
& \mbox{} + d_{13}(\t)\beta^{13}\geq 0\end{array}
\end{equation}

The inequalities (13) imply (14) because they imply that each of the
bracketed terms in (14) are non-negative for $0<\beta<t_1$.
Procedure 2 can be adapted (by changing the global variable {\sf c}
to {\sf d}) so that it can be applied to each of the trigonometric
coefficients $d_i(\t)$ to generate new polynomials $e_i(x)$.  Then,
Procedure 3 can be applied to each of the (transformed) linear
combinations specified in (13) to verify (14) and thus, complete the
case $0<\beta<t_1$.\\

\noindent{\bf Remarks.}\\

1.  We have verified the above constructions alternately using
Mathematica for the computer algebra component of the construction.

2.  The above process can be applied to $F_3(\a,x)$ and $F_5(\a,x)$
to give relatively straight-forward proofs of two cases of Brannan's
conjecture (6), specifically, the cases $2n+1=3$ and $2n+1=5$.  In
the latter case, the proof subdivides the interval $0<\a<1$ into two
cases $0<\a\leq 2/5$ and $2/5<\a<1$.  The argument here is
substantially simpler than Milcetich's proof.

3.  This technique for verifying Brannan's conjecture (6) for the
case $2n+1=7$ can be applied to an alternate, but closely related
coefficient inequality.  If in the series representation for
$F_N(\a,x)$ in (8) the summation is extended to infinity, that is, if
we write
\begin{eqnarray}
\nonumber
& F_N(\a,x)= & \sum^\i_{k=0} \sum^k_{j=0} \frac{(-\a)_j (-\a)_{k-j} 
(-1)^k \d_j \d_{k-j}}{k!(k-j)!}-\\[1ex]
\label{15}
&& \sum^\i_{k=0} \sum^k_{j=0} \frac{(-\a)_j (-\a)_{k-j} (-1)^k \d_j 
\d_{k-j} \cos((2*j-k)\t)}{k!(k-j)!}
\end{eqnarray}
and where again
$$
\d_j=\left\{ \begin{array}{ll}1&0\leq j\leq N\\[1ex] 
0&N+1\leq j\end{array}\right. ,
$$
then one can define the partial sums
\begin{eqnarray*}
& F^m_N(\a,x)= & \sum^m_{k=0}\sum^k_{j=0}\frac{(-\a)_j(-\a)_{k-j}
(-1)^k\d_j\d_{k-j}}{k!(k-j)!}-\\[1ex]
&& \sum^m_{k=0}\sum^k_{j=0}\frac{(-\a)_j(-\a)_{k-j}(-1)^k\d_j\d_{k-j}
\cos((2*j-k)\t)}{k!(k-j)!}
\end{eqnarray*}
Wheeler [Wh] considered the partial sums $F^m_N(\a,x)$.  He showed
that these partial sums have many properties which are analogous to
the coefficient sums $F_N(\a,x)$.  Specifically, he showed there that
for each $\a$, $0<\a<1$, there exists an $m_\a$ such that
$$
\max_{|x|=1}F^{2*m}_N(\a,x)<0
$$
for $m>m_\a$.  Furthermore, he devised the computer algebra technique
described above, and applied it to show that for $m=$ 1, 3, 5 and 7, 
$$
F^m_N(\a,x)\geq 0
$$
for $0<\a<1$.

\begin{thebibliography}{99}
\bibitem[AhFr]{1} D.\ Aharonov and S.\ Friedland, {\em On an
inequality connected with the coefficient conjecture for functions of
bounded boundary rotation}, Ann.\ Acad.\ Sci.\ Fenn.\ Ser.\ A 1
Math.\ {\bf 542} (1972), 14 pp.
\bibitem[Br1]{2} D.A.\ Brannan, {\em On functions of bounded boundary
rotation II}, Bull.\ London Math Soc.\ {\bf 1} (1969), 321--322.
\bibitem[Br2]{3} D.A.\ Brannan, {\em On coefficient problems for certain
power series}, Symposium on Complex Analysis, Canterbury, (1973)
eds.\ J.G.\ Clunie \& W.K.\ Hayman.  London Math.\ Soc.\ Lecture Note
Series {\bf 12} (1974), 17-27.
\bibitem[BrClKi]{4} D.A.\ Brannan, J.G.\ Clunie and W.E.\ Kirwan,
{\em On the coefficient problem for functions of bounded rotation},
Ann.\ Acad.\ Sci.\ Fenn.\ Ser.\ A 1 Math.\ {\bf 543} (1973), 18 pp.
\bibitem[BrHe]{5} G.\ Brown and E.\ Hewitt, {\em A class of positive
trigonometric sums}, Math.\ Ann.\ {\bf 268} (1984), 91--122.
\bibitem[Le]{6} O.\ Lehto, {\em On the distortion of conformal
mappings with bounded boundary rotation}, Ann.\ Acad.\ Sci.\ Fenn.\
Ser.\ A 1 Math.\ Phys.\ {\bf 124} (1952), 14 pp.
\bibitem[LoTa]{7} H.\ Lonka and H.\ Tammi, {\em On the use of
step-functions in extremal problems of the class with bounded
boundary rotation}, Ann.\ Acad.\ Sci.\ Fenn.\ Ser.\ A 1 {\bf 418}
(1968), 18 pp.
\bibitem[Mi]{8} J.G.\ Milcetich, {\em On a coefficient conjecture of
Brannan}, J.\ Math.\ Anal.\ Appl.\ {\bf 139} (1989), 515-522.
\bibitem[Pa]{9} V.\ Paatero, {\em \"{U}ber die konforme Abbildung 
von Gebieten deren R\"{a}nder von beschr\"{a}nkter Drehung sind},
Ann.\ Acad.\ Sci.\ Fenn.\ Ser.\ A {\bf 33} (1931), 77 pp.
\bibitem[ScTa]{10} M.\ Schiffer and O.\ Tammi, {\em On the fourth
coefficient of univalent functions with bounded boundary rotation},
Ann.\ Acad.\ Sci.\ Fenn.\ Ser.\ A {\bf 396} (1967), 26 pp.
\bibitem[Wh]{11} W.\ Wheeler, {\em Properties of partial sums of
certain special functions in complex variables}, Thesis, Texas Tech
University, 1995.
\end{thebibliography}

\noindent Department of Mathematics, Texas Tech University, Lubbock,
TX 79409\\
\noindent E-mail: barnard@math.ttu.edu, pearce@math.ttu.edu, 
bwheeler@math.ttu.edu.

\end{document}